\subsection{Signal analysis (H1--H3)}
\label{sec:results-signals}

We report whether each \emph{signal representation} predicts oracle $r(q)$ on calib before any \textbf{binary routing} policy claims.
H1--H3 are \textbf{parallel representation tests}: each asks whether that layer's feature vector helps predict $r(q)$, not whether one layer beats another and not whether $\phi{+}\psi$ combination defines H3 (H3 = $\chi$ alone; combinations support Stage~7 selection).

\subsubsection{Representation evaluation (6.1)}
\label{sec:results-signals-representation}

For each representation on $R_c$, we fit a logistic evaluation model (5-fold stratified CV) and report mean CV AUROC and AUPRC vs.\ $r(q)$, with per-fold scores stored in \texttt{representation\_tests.json} for stability analysis.
Representations: \textbf{H1a} structural, \textbf{H1b} ambiguity, \textbf{H1c} embedding geometry, \textbf{H1} combined $\phi$, \textbf{H2} $\psi(M_{\mathrm{lo}})$, \textbf{H3} $\chi$.
Table~\ref{tab:signal-analysis} is the primary evidence for H1--H3.
\textsc{[Values from \texttt{representation\_tests.json}.]}

\subsubsection{Feature diagnostics (6.2)}
\label{sec:results-signals-diagnostics}

Per-feature Spearman $\rho$, AUROC, AUPRC, and class-conditional means ($r{=}0$ vs.\ $r{=}1$) from \texttt{univariate\_\{query,response,cross\}.jsonl} support interpretation of \emph{why} a representation may work; they are not the primary hypothesis tests.
\textsc{[Per-feature appendix tables TBD.]}

\begin{table}[t]
  \centering
  \small
  \begin{tabular}{@{}lcc@{}}
    \toprule
    \textbf{Representation} & \textbf{Joint CV AUROC} & \textbf{Joint CV AUPRC} \\
    \midrule
    H1: structural ($\phi_{\text{structural}}$, H1a) & \textsc{TBD} & \textsc{TBD} \\
    H1: ambiguity ($\phi_{\text{ambiguity}}$, H1b) & \textsc{TBD} & \textsc{TBD} \\
    H1: embedding geometry ($\phi_{\text{geometry}}$, H1c) & \textsc{TBD} & \textsc{TBD} \\
    H1: all $\phi$ combined (H1) & \textsc{TBD} & \textsc{TBD} \\
    H2: $\psi(M_{\mathrm{lo}})$ & \textsc{TBD} & \textsc{TBD} \\
    H3: $\chi$ & \textsc{TBD} & \textsc{TBD} \\
  \bottomrule
  \end{tabular}
  \caption{Representation-level signal analysis on calib vs.\ oracle $r(q)$. Each row tests whether that \emph{representation} (feature vector) predicts $r(q)$ via 5-fold stratified CV logistic regression---not router training. Per-scalar diagnostics appear in the appendix; feature ranking and redundancy analysis are Stage~7 (Section~\ref{sec:signal-selection}).}
  \label{tab:signal-analysis}
\end{table}


\paragraph{Summary.}
Table~\ref{tab:signal-analysis} asks whether each representation---not each individual scalar---contains routing information.
A positive H2 with a weak H1 is valid---model-dependent uncertainty may carry routing information without query-side cues.
Feature ranking, redundancy pruning, and frozen $x(q)$ are Stage~7 (Section~\ref{sec:signal-selection}).
